\documentclass[a4paper,11pt]{amsart}

\usepackage[all]{xy}
\usepackage{amscd}
\usepackage{amsfonts}
\usepackage{a4wide}
\usepackage{amsmath}
\usepackage{amssymb}
\usepackage{mathtools}
\usepackage{natbib}
\usepackage{fullpage}
\usepackage[applemac]{inputenc}

\usepackage{graphicx}

\addtolength{\voffset}{-2cm}
\addtolength{\textheight}{3cm}


%\usepackage[colorlinks=true]{hyperref}

\usepackage[pdftex,backref,hyperindex,breaklinks,colorlinks,citecolor=blue]{hyperref}
%index (indice de palabras)


%letras griegas

\newcommand{\sg}{\sigma}
\newcommand{\eps}{\varepsilon}
\newcommand{\ph}{\varphi}


%\newcommand{\integral}[2]{\int{#1}\,\mathrm{d}{#2}}
%\integral{funcion}{medida}

%cuerpos y anillos basicos

\newcommand{\Fp}{\mathbb{F}_p}
\newcommand{\Fpbar}{\overline{\mathbb{F}}_p}
\newcommand{\ce}{\mathbb{C}}
\newcommand{\erre}{\mathbb{R}}
\newcommand{\z}{\mathbb{Z}}
\newcommand{\ene}{\mathbb{N}}
\newcommand{\q}{\mathbb{Q}}


%superficies de Riemann basicas

\newcommand{\h}{\mathbb{H}}
\newcommand{\p}{\mathbb{P}^1}
\newcommand{\disco}{\mathbb{D}}


%objetos modulares

\newcommand{\XC}{X_0}
\newcommand{\XQ}[1]{X_0 (#1)_\q}
\newcommand{\XZ}{\mathcal{X}_0}
\newcommand{\YC}{Y_0}
\newcommand{\YZ}{\mathcal{Y}_0}
\newcommand{\G}{\Gamma_0}
\newcommand{\T}{\hat{\mathbb{T}}}
\newcommand{\TN}{\hat{\mathbb{T}}_N}
\newcommand{\TNH}{\hat{\mathbb{T}}_{N,H}}
\newcommand{\TNQ}{\mathbb{T}_N^0}
\newcommand{\zloc}[1]{\mathbb{Z}[\frac{1}{#1}]}
\newcommand{\Ch}{\widehat{CH} (N)}
\newcommand{\ChR}{\widehat{CH}_\erre (N)}
\newcommand{\Chnum}{\widehat{CH}_\erre^{num} (N)}
\newcommand{\Chu}{\widehat{CH} (1)}
\newcommand{\infinito}{\hat{D}_\infty}
\newcommand{\JNQ}{J_0 (N)_\q}
\newcommand{\JNR}{J_0 (N)_{\erre}}
\newcommand{\JN}{J_0 (N)}
\newcommand{\du}{\frac{dxdy}{y^2}}


%geometria 

\newcommand{\Sch}{\mathcal{D}}
\newcommand{\X}{\mathcal{X}}
\newcommand{\Y}{\mathcal{Y}}
\newcommand{\CH}{\widehat{CH^1}}
\newcommand{\dv}{\textrm{div}}
\newcommand{\Spec}{\textrm{Spec}}
\newcommand{\Lb}{\mathcal{L}}
\newcommand{\Lbhat}{\widehat{\mathcal{L}}}
\newcommand{\Aan}{\mathbb{A}^{2,an}}
\newcommand{\Oo}{\mathcal{O}}



%pedanterias

\newcommand{\der}{\mathclose>}
\newcommand{\izq}{\mathopen<}

%miscelaneos

\newcommand{\comb}[2]{\left(\begin{array}{c} 
#1 \\
#2 
\end{array}\right)}
\newcommand{\wbar}{\hat{\omega}}
\newcommand{\dosderi}{\frac{\partial^2}{\partial z \partial \bar{z}}}
\newcommand{\deri}{\frac{\partial}{\partial z}}
\newcommand{\deriw}{\frac{\partial}{\partial w}}
\newcommand{\dz}{ dz \wedge d \bar{z}}
\newcommand{\epsi}{\epsilon \rightarrow 0}
\newcommand{\Luno}{L^2_1}
\newcommand{\Nolu}{L^2_{-1}}
\newcommand{\pruf}{\noindent \textbf{Proof}: }
\newcommand{\M}{\mathcal{M}}
\newcommand{\MFS}{\M^{1}}
\newcommand{\LL}{\overline{\mathcal{L}}}
\newcommand{\Qbar}{\overline{\mathbb{Q}}}
\newcommand{\MZ}{\mathcal{P}^{\overline{\z}}_{\operatorname{log}}}
\newcommand{\dd}{\mathrm{d}}
\newcommand{\ha}[1]{\widehat{#1}} 
\renewcommand{\and}{\quad \text{and} \quad}

%numeracion
\renewcommand{\theequation}{\thesection.\arabic{equation}}


%teoremas, lemas, etc



\newtheorem{thm}{Theorem}[section]
\newtheorem{lema}[thm]{Lemma}
\newtheorem{prop}[thm]{Proposition}
\newtheorem{prop-def}[thm]{Proposition-Definition}
\newtheorem{cor}[thm]{Corollary}
\newtheorem{conj}[thm]{Conjecture}
\theoremstyle{definition}
\newtheorem{rem}[thm]{Remark}
\newtheorem{rems}[thm]{Remarks}
\newtheorem{ex}[thm]{Example}
\newtheorem{definicion}[thm]{Definition}

\title{Around the Gauss point in dimension 2}

\author{Eyal Goren}

\author{Sebasti\'an Herrero}

\author{Ricardo Menares}
\address{Pontificia Universidad Cat\'olica de Chile}
\email{rmenares.v@gmail.com}


\begin{document}

\maketitle

\tableofcontents

\section{The Berkovich affine plane}

Let $p$ be a prime number. We denote by $|\cdot|$ the $p$-adic norm, extended from the standard norm on $\q_p$ to $\ce_p$. Set $\Oo_p:=\{ z \in \ce_p : |z|\leq 1\}$ and fix once and for all a reduction map $\Oo_p\rightarrow \Fpbar$. For any $z\in \Oo_p$, we denote by $\tilde{z}$ its image in $\Fpbar$.

\begin{definicion}
The Berkovich affine plane $\Aan$ over $\ce_p$ is the analytic spectrum of $\ce_p[S,T]$. As a set, $\Aan$ is the collection of all multiplicative seminorms on $\ce_p[S,T]$ extending the norm on $\ce_p$. The topology in $\Aan$ is the weakest topology such that for all $f\in \ce_p[S,T]$, the map $\iota_f : \Aan \rightarrow \erre_{\geq 0}$ given by $\iota_f(x)=|f|_x$, is continuous.

The Gauss point $\xi \in \Aan$ is the norm defined as $$\left| \sum_{i,j}a_{i,j}S^iT^j\right|_\xi := \max\{|a_{i,j}|\}.$$ 
\end{definicion}


\begin{lema}\label{vecindades}
For any $R>1$ and $f\in \ce_p[S,T]$, set $$V(f;R)=\{ x \in \Aan : R^{-1}<|f(S,T)|_x<R\}$$
A system of neighbourhoods of the Gauss point is $$\left\{V(f;R) : \quad f \in \ce_p[S,T], \quad |f|_\xi=1, \quad  R> 1\right\}.$$
\end{lema}

\begin{proof}
The topology of $\Aan$ is generated by sets formed by choosing an open bounded interval $I\subseteq [0,\infty)$, an element $f\in \ce_p[S,T]$ and setting

$$U(f,I)=\{ x \in \Aan : |f|_x \in I \}.$$ 

Note that for any $\lambda \in \ce_p^*$, we have that $$U(\lambda f,I)=U(f,|\lambda|^{-1}I).$$ Hence, we can restrict our attention to the collection
$$\{U(f,I) : |f|_\xi=1, \quad I\subset [0,\infty) \textrm{ bounded open interval}\}.$$ 

Such a set contains the Gauss point if and only if $1\in I$. Hence, we can restrict further to intervals of the form $I=(R^{-1},R)$ with $R>1$, as desired. 
\end{proof}


\begin{lema}\label{criterion}
Let $(D_n)\subseteq \ce_p^2$ be a sequence of effective 0-cycles such that $\deg D_n\geq 1$ for all $n$. Then, the following are equivalent:
\begin{enumerate}
\item[i)] $\overline{\delta}_{D_n} $ weakly converges to $\delta_\xi$
\item[ii)] For every $R>1$ and $f\in \ce_p[S,T]$ with $|f|_\xi=1$, we have that $$\lim_{n\rightarrow \infty} \frac{\deg \left(D_n\cap V(f;R)\right)}{\deg D_n} =1.$$ 
\end{enumerate}
\end{lema}

\begin{proof}
Since the measures in consideration have finite support, weak convergence is equivalent to the statement that for every neighbourhood $U$ of $\xi$, it holds $$\lim_{n\rightarrow \infty} \frac{\deg (D_n\cap U)}{\deg D_n} =1.$$ Hence, the assertion follows from Lemma \ref{vecindades}.
\end{proof}

\begin{lema}\label{criterionmodp}
Let $(D_n)\subseteq \Oo_p^2$ be a sequence of effective $0$-cycles such that for any curve $C\subseteq \mathbb{A}^2_{\Fp}$, it holds $$\lim_{n\rightarrow \infty} \frac{\deg\left(\tilde{D}_n\cap C(\Fpbar)\right)}{\deg D_n} =0.$$ Then, $\overline{\delta}_{D_n}$ weakly converges to $\delta_\xi$.
\end{lema}

\section{Hecke correspondences on the bad locus}
Let $Y$ be the modular curve. Let $\ell$ be a prime number different from $p$ and consider the Hecke correspondence $T_\ell$ on $Y\times Y$. In modular terms, we have that for any pair of elliptic curves $(E_1,E_2)$ it holds $$T_\ell(E_1,E_2)=\sum_{C_1\leq E_1 \atop |C_1|=\ell } \sum_{C_2\leq E_2 \atop |C_2|=\ell } (E_1/C_1,E_2/C_2).$$

In particular, $\deg T_\ell = (\ell+1)^2$. 

Identify $Y(\ce_p)$ with $\ce_p$ through the $j$-invariant. Then, the analytification of $Y_{\ce_p}\times Y_{\ce_p}$ is isomorphic to $\Aan$. Let $Y_{bad}$ be set of points in $Y(\ce_p)$ parametrizing elliptic curves with bad multiplicative reduction. Then,  Hecke correspondences preserve $Y_{bad}\times Y_{bad}$. 

\begin{prop} \label{badlocus}
Let $P\in Y_{bad}\times Y_{bad}$. Then, we have that $\overline{\delta}_{T_\ell(P)} \rightarrow \delta_{\xi}$ weakly as $\ell \rightarrow \infty$ among prime numbers. 
\end{prop}

Before proving Proposition \ref{badlocus} we state an elementary fact. 

\begin{lema}\label{elementary}
Let $m,A$ be positive integers and let $q\in \ce_p$ with $|q|=1$. Then, there exist $t\geq 1$ such that for all integers $\ell\geq 1$ with $p\nmid \ell$ it holds
\begin{equation}\label{goal2}
|\{ \gamma \in \ce_p : \gamma^{m\ell}=q \textrm{ and } [\Fp(\tilde{\gamma}):\Fp]\leq A\}| \leq t.
\end{equation}
\end{lema}

\begin{proof}\footnote{Perhaps one can simplify the exposition a bit. First one notes the solutions are in $\Oo_p$. I think that it logically follows -- just from the statement -- that the lemma is true for a general $m$ if and only if it is true for $m$ of the form $p^r$. Then, as the map $\gamma \rightarrow \gamma^{p^r}$ is finite to 1, and $\mathbb{F}_p(\tilde \gamma) = \mathbb{F}(\widetilde{\gamma^{p^r}})$, one may assume $m=p^r=1$. Then one gives your proof, but it might be easier to follow now. \textcolor{blue}{RM: you are right. If we end publishing this particular piece of the problem I will do the modification you suggest} (A side remark is that one usually uses $q$ for a power of $p$, so it might be better to replace $q$ by $u$ or $z$, say.\textcolor{blue}{RM: yes, but here $q$ is a Tate parameter})}
Set $n:=\ell m$. All roots of $X^{n} = q$ lie in $\Oo_p$ because $|q|=1$. Choose one root $\gamma_0$. Then, the set of all roots is $\{\gamma_0\zeta : \zeta\in\mu_n\}$, where $\mu_n$ stands for the group of $n$-th roots of unity in $\ce_p$.

Write $m = p^r m_0$ with $p\nmid m_0$ and set $n' := \ell m_0$. The restriction of the reduction map $\rho:\mu_n \to \overline{\mathbb{F}}_p^\times$ is a homomorphism. Since $p\nmid \ell$, its kernel consists exactly of the $p^r$-th roots of unity, and its image is $\mu_{n'}$. Consequently, for a fixed $y\in\overline{\mathbb{F}}_p^\times$, the equation $\rho(\zeta)=y$ has either $0$ or $p^r$ solutions $\zeta\in\mu_n$.

Let
\[
R = \{y\in\overline{\mathbb{F}}_p^\times : y^n = \widetilde{q}\}, \quad S = \bigl\{ y\in\overline{\mathbb{F}}_p^\times : [\mathbb{F}_p(y):\mathbb{F}_p]\le A \bigr\},\]
\[ 
T=\{ \gamma \in \ce_p : \gamma^{n}=q \textrm{ and } [\Fp(\tilde{\gamma}):\Fp]\leq A\}.
\]
For any $\gamma \in T$, its reduction $\widetilde{\gamma}$ lies in $S\cap R$. The set $R$ is precisely the coset $\widetilde{\gamma}_0\mu_{n'}$ and each $y\in S\cap R$ lifts to exactly $p^r$ distinct roots  of $X^n=q$. Therefore, 
\[
|T|\leq p^r \cdot |S\cap R| \le p^r |S|, 
\]
showing \eqref{goal2} with  $t = p^r |S|$.
\end{proof}

\begin{proof}[Proof of Proposition \ref{badlocus}]
The Tate uniformization of elliptic curves with bad multiplicative reduction furnishes an identification of $Y_{bad}$ with $\{ q \in \ce_p : |q|<1\}$. In terms of the parameter $q$, we have the following description of $T_\ell$. Let $(q_1,q_2)$ be the parameters corresponding to  $(E_1,E_2)$. Let $$Q_\ell=\{ (\eta_1,\eta_2) : \eta_1^\ell=q_1, \eta_2^\ell=q_2\}.$$ Then, 
\begin{equation}\label{formula}
T_\ell(q_1,q_2)=\sum_{(\eta_1,\eta_2)\in Q_\ell} (\eta_1,\eta_2)+\sum_{\eta_1^\ell=q_1}(\eta_1,q_2^\ell)+\sum_{\eta_2^\ell=q_2}(q_1^\ell, \eta_2)+(q_1^\ell,q_2^\ell).
\end{equation}
\footnote{Is this true?? Don't you need terms like $\sum_{\eta_1\in Q_\ell} (\eta_1, q_2^\ell)$? The sentence "We remark that..." below suggest that you simply forgot the $\sum$ sign.\textcolor{blue}{RM: yes, the writing of the formula was ambiguous. I corrected it.}}

In view of Lemma \ref{criterion}, we need to show that for every $f\in \ce_p[S,T]$ with $|f|_\xi=1$ and every $R>1$, we have that 
\begin{equation}\label{goal}
\lim_{\ell \rightarrow \infty}\frac{\deg (T_\ell(q_1,q_2)\cap V(f;R))}{(\ell+1)^2}=1.
\end{equation}

We remark that in \eqref{formula}, the terms $(\eta_1,q_2^\ell)+(q_1^\ell, \eta_2)$ and $(q_1^\ell,q_2^\ell)$ account for only $2\ell+1$ points, so they have no influence in the computation of the limit in \eqref{goal}. We will focus on the terms of the form $(\eta_1,\eta_2)$. 

Write $f(S,T)=\sum_{i,j}a_{i,j}S^iT^j$. Let\footnote{Using $T$ in the following displayed line is really bad notation\textcolor{blue}{RM: oh yes}} $$I=\{(i,j) : |a_{i,j}|=1\}, \quad \kappa:=\max_{(i,j) \in I} |q_1|^i|q_2|^j, \quad C=\{(i,j) \in I :  |q_1|^i|q_2|^j=\kappa \}, \quad M=\max_{(i,j)\notin I}|a_{i,j}|.$$ 

Assume first that $C$ is a singleton, $C=\{(i_0,j_0)\}$. Then for every $(i,j)\notin I$, we have that $|a_{i,j}\eta_1^i\eta_2^j|<M$. For $(i,j)\in I$ with $(i,j)\neq (i_0,j_0)$, we have that $$|a_{i,j}\eta_1^i\eta_2^j|=|q_1|^{i/\ell}|q_2|^{j/\ell}<|q_1|^{i_0/\ell}|q_2|^{j_0/\ell} =|a_{i,j} \eta_1^{i_0}\eta_2^{j_0}|.$$

Since $M<1$ and $|a_{i,j} \eta_1^{i_0}\eta_2^{j_0}|=|q_1|^{i_0/\ell}|q_2|^{j_0/\ell}$ tends to 1 as $\ell \rightarrow \infty$, we have that for $\ell$ big enough it holds\footnote{You write "we have that for $\ell$ big enough it holds.." but then write a limit... \textcolor{blue}{RM: I agree it's too informal, I hope now it's better}} $$|f(\eta_1,\eta_2)|=|q_1|^{i_0/\ell}|q_2|^{j_0/\ell}$$ Hence, $|f(\eta_1,\eta_2)| \rightarrow 1 $  as $\ell \rightarrow \infty$, thus establishing \eqref{goal} in the case where $C$ is a singleton.
\footnote{Is this where you want to finish the discussion, or do you want already to conclude (\ref{goal}) at this point?? \textcolor{blue}{RM: this proves \eqref{goal} in this case, I hope now it's more clear}}

Now assume $|C|\geq 2$. Consider the auxiliary polynomial $$g(S,T)=\sum_{i,j \in C}a_{i,j}S^iT^j.$$ 

{\bf Claim}. There are positive integers $\alpha, \beta, k$ such that
\begin{enumerate}
\item[i)] for any $(i,j)\in C$ it holds $k=i\alpha + j\beta$.
\item[ii)] $\left|\frac{q_2^\alpha}{q_1^\beta}\right|=1$.
\end{enumerate}

Indeed, write $|q_i|=p^{-r_i}$ with $r_1,r_2\in \q_{>0}$. Take distinct elements $(i,j),(i',j')\in C$. Then, from $|q_1|^i|q_2|^j=|q_1|^{i'}|q_2|^{j'}$ we deduce $ir_1+jr_2=i'r_1+j'r_2$. Writing $r_i=\frac{a_i}{b_i}$ with $a_i,b_i\geq 1$, we see that  $\alpha=a_1b_2, \beta=a_2b_1$ satisfy the conditions in the claim.\\



 Now set $h(T):=\sum_{(i,j)\in C} a_{i,j}T^{j\beta}$ and for every $(\eta_1,\eta_2)\in Q_\ell$, choose once and for all $(\xi_1,\xi_2)\in \ce_p^2$ such that $\xi_1^{\alpha}=\eta_{1}, \xi_2^\beta=\eta_{2}$. Then, we see that 
 $$g(\eta_1,\eta_2)=g(\xi_1^\alpha,\xi_2^\beta)=\sum_{(i,j)\in C}a_{i,j}\xi_1^{i\alpha}\xi_2^{j\beta}=\xi_1^kh\left(\frac{\xi_2}{\xi_1}\right).$$


Now there are a finite number of elements $a_1,\ldots, a_r\in \ce_p$ such that $|h(z)|<1$ if and only if $z\in \cup_{i=1}^rD^0(a_i,1).$ So $\left |h\left(\frac{ \xi_2}{\xi_1}\right)\right|<1$ if and only if there is some $i\in \{1,\ldots,r\}$ such that $$\left|\frac{ \xi_2}{\xi_1}-a_i\right|<1.$$

Since $$\left(\frac{\xi_2}{\xi_1}\right)^{\ell \alpha \beta}=\frac{q_2^{\alpha}}{q_1^\beta}=:q$$ and $|q|=1$ due to the choice of $\alpha, \beta$, using Lemma \ref{elementary} we see that there is $t\geq 1$ such that  for each $i=1,2,\ldots,r$ and all $\ell \neq p$, it holds
$$\left|\left\{(\eta_1,\eta_2)\in Q_\ell : \left|\frac{ \xi_2}{\xi_1}-a_i\right|<1\right\}\right|\leq \ell t.$$
In conclusion, for all $\ell\neq p$ there exists a subset $R_\ell \subset Q_\ell$ such that 
\begin{enumerate}
\item[i)] $|R_\ell|\geq |Q_\ell|-\ell rt=\ell^2-\ell rt$
\item[ii)] for all $(\eta_1,\eta_2)\in R_\ell$, it holds $|g(\eta_1,\eta_2)|=|\xi_1|^k=|q_1|^{\frac{k}{\ell \alpha}}$
\end{enumerate}

Since $$|f(\eta_1,\eta_2)-g(\eta_1,\eta_2)|<M<1$$ and   $|q_1|^{\frac{k}{\ell \alpha}}$ tends to 1 as $\ell \rightarrow \infty$, we conclude that for $\ell$ big enough

$$\frac{\deg (T_\ell(q_1,q_2)\cap V(f;R))}{(\ell+1)^2}\geq \frac{\ell^2-r\ell t}{(\ell+1)^2},$$
and so \eqref{goal} holds, as desired.
\end{proof}


\iffalse

\section{Shilov boundary of Arithmetic varieties}

\begin{definicion}
Let $A$ be a Banach algebra. A set $\Gamma \subset M(A)$ is a \emph{boundary} if $\Gamma$ is closed and every $a\in A$ reaches its maximum in $\Gamma$.

Using Zorn's Lemma, one can show that there are always minimal boundaries (with respect to $\subseteq$). If there exists a \emph{unique} minimal boundary, it will be called the \emph{Shilov boundary}.  
\end{definicion}

reduction map

\begin{thm}
Gauss points and generic points of irreducible components
\end{thm}

\begin{prop} (Mod p criterion)
Sufficient condition: convergence towards the generic point mod p implies convergence towards the Gauss point
\end{prop}

\fi

\bibliographystyle{alpha}



\bibliography{papers}





\end{document}


%%% Local Variables:
%%% mode: latex
%%% TeX-PDF-mode: t
%%% TeX-source-correlate-mode: t
%%% TeX-master: t
%%% End: 
